\newpage
\centeredheading{ЗАКЛЮЧЕНИЕ}
\addcontentsline{toc}{chapter}{ЗАКЛЮЧЕНИЕ}

К моменту окончания учебной практики все задачи, предусмотренные индивидуальным заданием, были успешно выполнены в полном объеме.
Участие в разработке прототипа системы AutoSOAR позволило мне получить уникальный инженерный и исследовательский опыт на стыке трех передовых направлений: безопасной разработки, наступательной безопасности и применения технологий искусственного интеллекта в ИБ\@.

В результате проделанной работы были достигнуты следующие ключевые результаты:
\begin{enumerate}
    \item С нуля организовано рабочее пространство проекта в GitHub Organization.
    Внедрена методология Agile, настроена Канбан-доска, определен workflow безопасной совместной разработки.
    \item Успешно устранен архитектурный Vendor Lock-in в модуле классификации инцидентов.
    Переход на универсальную фабрику обеспечил гибкость и независимость от конкретного провайдера LLM, а также позволил использовать локальные модели.
    \item Опираясь на матрицу MITRE ATT\&CK, разработаны сценарии эмуляции кибератак высокой сложности.
    Данные сценарии стали эталоном для генерации обучающей и тестовой телеметрии на киберполигоне.
\end{enumerate}

Практика продемонстрировала важность автоматизации процессов в современных центрах мониторинга и показала эффективность применения ИИ-агентов для первичного анализа инцидентов.

Полученные мной компетенции и опыт подкрепляют и делают осознанным мой выбор образовательной траектории «Безопасная разработка ПО».